\documentclass[11pt]{article}
\usepackage[margin=1in]{geometry}
\usepackage{amsmath,amssymb,amsthm,hyperref}
\newtheorem{theorem}{Theorem}
\newtheorem{proposition}[theorem]{Proposition}
\newtheorem{lemma}[theorem]{Lemma}
\newcommand{\MA}{\operatorname{MA}}
\newcommand{\DD}{\mathrm d'\mathrm d''}
\newcommand{\C}{\mathcal C_\theta}
\newcommand{\R}{\mathbb R}
\title{Multiscale branch tests and the critical tail of a bounded energy dual}
\author{Research working paper, continuation 384}
\date{9 October 2026}
\begin{document}
\maketitle
\begin{abstract}
We prove a precise implication within an explicitly specified bounded analytic
calculus. A single branch pencil has finite dual exactly at the weak tail
threshold $t^{-2/3}$. Finite positive combinations of branch cutoffs instead
have finite dual exactly when the tail belongs to $L^{3/2}(\mathrm dt)$.
A countable analytic atomic measure at the critical power passes every
single-pencil boundedness test but has infinite dual on the finite combinations.
The bounded calculus and quotient geometry are the unreviewed inputs of
continuation 383; their independent review is pending. This paper proves the
multiscale implication, not those foundations or the full original problem P16.7.
\end{abstract}

\section{Exact standing assumptions and the branch}
We use the explicit analytic surface of continuation 383, with
$K=\mathbb C((s))$, $|s|=e^{-1/8}$, $t=s^8$, and periods
$(t,2),(3,t)$ on the two-dimensional analytic torus. Its good, separated,
proper, boundaryless quotient geometry and nonalgebraization remain subject
to independent verification. On a deck sheet put $x_j=-\log|z_j|$,
$q=(x_1^2+x_2^2)/2$, and $\theta=\DD q$.
The bounded class $\C$ consists of continuous real functions $u$ for which
$q+u$ is, locally on ordinary opens, a locally uniform limit of ordinary
CLD smooth functions with positive semidefinite Hessian. We do not identify
this class with an intrinsic unbounded plurisubharmonic class.

Here are the precise calculus hypotheses used in this paper.
\begin{enumerate}
\item The class $\C$ is convex. Its locally defined bounded mixed products
are positive actual analytic Radon measures, symmetric and bilinear in the
Hessians, local on ordinary opens, continuous under uniform approximation
within the class, and have raw top mass two. Continuous coefficient
multiplication and the differentiated-product identity used below are valid.
\item Write $G(y)=y^2/2-y/2+1/12$ on $[0,1]$, periodically extended,
$u_0=-G(x_1)-G(x_2)$, and $A=\theta+\DD u_0$.
The bounded energy and reference identities are
\[
 E(u)=\frac13\int_X u\bigl(\theta^2+\theta\wedge(\theta+\DD u)
                         +\MA(u)\bigr),\qquad E(u_0)=-\frac16,
\]
\[
 A^2=2\delta_{i(0)},\quad
 \theta^2=2i_*\mathrm dx,\quad
 \theta\wedge A=i_*(\mathrm dy\otimes\delta_0+
                                      \delta_0\otimes\mathrm dy).
\]
Here $i$ is the ordinary torus Gauss section, $\mathrm dx$ is probability
Haar measure on $\R^2/\mathbb Z^2$, and the entire fibre
$F=\tau^{-1}(0)$ has $u_0=-1/6$.
\item On an ordinary translated identity torus chart with coordinates
$w_1,w_2$, the positive top measure of a smooth convex weight $f(y)$,
$y_j=\log|w_j|$, is supported on its actual Gauss section and equals
$2\det(\operatorname{Hess} f)\,\mathrm dy$ there. The section is continuous
against arbitrary continuous analytic tests. This is the degree-one,
saturated-lattice calibration with the exterior-product factor $2!$.
\end{enumerate}
These are explicitly conditional inputs, not newly accepted correctness
grants. Continuation 383 supplies complete proposed proofs and its source
ledger for them. The numerical mass is never replaced by probability mass.

Choose the ordinary sheet $V=\{|x_j|<1/4\}$ over $F$. For every real $N\geq1$,
put $r=e^{-N}$ and, on $V$,
\[
 \ell_j=\log\max(r,|z_j-1|),\quad
 b_j=\log\max(1,|z_j|),\quad
 D_N=\max(\ell_1-b_1,\ell_2-b_2),\qquad U_N=u_0+D_N.
\]
The functions $D_N$ vanish off $F$ and extend by zero outside $V$.
Their negative locus is the ordinary open
\[
 W=\{|z_1-1|<1,\ |z_2-1|<1\}\subset F.
\]
With the extended convention $-\log0=+\infty$, define the Borel depth
\[
 T(p)=\min\{-\log|z_1(p)-1|,-\log|z_2(p)-1|\}\quad(p\in W),
 \qquad T=0\quad(p\notin W).
\]
Then $-D_N=\min(N,T)$ everywhere. In particular $T(y)=+\infty$ at the
rigid point $y=(1,1)$. The actual translated bidisc Gauss point $\xi_N$
has $|z_j-1|=e^{-N}$ and $T(\xi_N)=N$. Real radii require no assertion
that every such bidisc is strict.

The local identity
\[
 q+U_N=\max(\ell_1+b_2,b_1+\ell_2)
                    -\tfrac12\log|z_1z_2|-\tfrac16
\]
shows $U_N\in\C$: on nonzero-coordinate charts all participating logarithms
are ordinary moments, the displayed maxima are convex, and the last term
is pluriharmonic. Near a zero of $z_j-1$ its inactive capped logarithm can
be omitted. Smooth convex regularizations then supply the required local
approximants. Also $D_N\theta=D_N A=0$ since both currents vanish on $W$
and $D_N$ vanishes on its complement. These claims use precisely the
bounded calculus assumptions above.

\section{The measure and energy of every finite combination}
Let $1\leq N_1<\cdots<N_k<\infty$, $a_j>0$, and
$\sum_j a_j\leq1$. Zero coefficients and repeated cutoffs may be removed
or combined. Put
\[
 u_h=u_0+\sum_{j=1}^k a_jD_{N_j},\quad
 h(t)=\sum_j a_j\mathbf1_{[0,N_j)}(t),\quad
 g(t)=\int_0^t h(v)\,\mathrm dv=\sum_j a_j\min(t,N_j).
\]
Thus $h$ is nonnegative, nonincreasing, compactly supported, at most one,
and constant on $[0,1)$. Conversely every finite step function with these
properties has this form, by taking its nonnegative downward jumps.
Convexity gives $u_h\in\C$, and $u_h-u_0=-g(T)$.

\begin{proposition}\label{prop:energy}
Under the standing bounded calculus assumptions,
\[
 \MA(u_h)|_W=\sum_{j=1}^k
       \bigl(h(N_j-)^2-h(N_j+)^2\bigr)\delta_{\xi_{N_j}},
 \qquad E(u_h)=-\frac16-\frac13\int_0^\infty h(t)^3\,\mathrm dt.
\]
The remaining top mass outside $W$ is not identified or prescribed.
\end{proposition}
\begin{proof}
The zero profile gives $u_h=u_0$ and both identities immediately.
For the following gradient argument assume $H=\sum_j a_j>0$.
On $W$ the reference weight $q+u_0$ is constant, and the local weight is,
up to its constant,
\[
 f(y_1,y_2)=\sum_j a_j\max(-N_j,y_1,y_2).
\]
Here is a direct calibrated mass calculation, including analytic tests.
Smooth it by
\[
 f_\epsilon(y)=\sum_j a_j\epsilon
        \log(e^{-N_j/\epsilon}+e^{y_1/\epsilon}+e^{y_2/\epsilon}).
\]
Its Hessian is positive definite. If $H=\sum_j a_j$, its gradient maps
$\R^2$ bijectively onto the interior of
$\Delta_H=\{p_1,p_2\geq0,\ p_1+p_2\leq H\}$.
Indeed, the strict convexity gives injectivity. For each interior $p$,
$f_\epsilon(y)-p\cdot y$ is coercive because its recession function is
$H\max(0,y_1,y_2)-p\cdot y>0$ for every nonzero recession direction;
its minimum supplies the inverse gradient point.

The convergence $f_\epsilon\to f$ is uniform with error at most
$\epsilon H\log3$. For $p_1,p_2>0$ and
\[
 h(N_j+)<p_1+p_2<h(N_j-),
\]
the unique minimizer of $f(y)-p\cdot y$ is $(-N_j,-N_j)$.
To see uniqueness, for fixed $\max(y_1,y_2)$ the positive coordinates
of $p$ make both $y_j$ attain that maximum. The remaining scalar
convex piecewise affine function has its unique minimum at the stated
break because its slopes straddle $p_1+p_2$ strictly.
Uniform convergence and coercivity therefore make the inverse gradient
points converge to this vertex. The exceptional boundaries in $p$ have
two-dimensional Lebesgue measure zero.

Change variables by the gradient in $2\det(\operatorname{Hess}f_\epsilon)$.
For bounded continuous tests dominated convergence gives the mass
$h(N_j-)^2-h(N_j+)^2$ at that vertex: the intervening trapezoid has
area one half of this difference. Composing a continuous analytic test
with the actual translated Gauss section proves the same statement
against analytic tests, not just their tropical restrictions.
This calculation on the invertible-$w$ locus determines the top measure
there by local product continuity. At other points the unsmoothed weight
is locally affine or depends, after an affine subtraction, on one moment,
so its square vanishes. Near a zero of one $w_j$ all its finitely many
capped logarithms are inactive, leaving at most one active moment;
near the simultaneous zero the weight is constant. Thus there are no
additional points of top mass inside $W$.

Let $D=u_h-u_0=-g(T)$. We have $D\theta=DA=0$.
The differentiated-product identity implies
$A\wedge\DD D=\DD(DA)=0$ and
$\theta\wedge(\theta+\DD u_h)=\theta\wedge A$.
By ordinary locality the top measure is supported on $F$, where
$u_0=-1/6$, and its total mass is two. The first two energy integrals
are consequently $0$ and $-1/6$; the third is
$-1/3+\int D\,\MA(u_h)$. The displayed atom formula gives
\[
 -\int D\,\MA(u_h)=\sum_j g(N_j)
                \bigl(h(N_j-)^2-h(N_j+)^2\bigr)
                         =\int_0^\infty h(t)^3\,\mathrm dt.
\]
The last equality is ordinary integration by parts for a finite step
function: its square has the indicated downward jumps, $g(0)=0$,
$h$ vanishes at infinity, and $g'=h$ almost everywhere.
Substituting proves the energy identity.
\end{proof}

\section{Two exact restricted duals}
Let $\mu$ be any positive analytic Radon measure of raw mass two, and put
\[
 J_0=E(u_0)-\int u_0\,\mathrm d\mu,\quad
 F_\mu(t)=\mu\{T>t\},\quad
 L_N=\int\min(N,T)\,\mathrm d\mu=\int_0^N F_\mu(t)\,\mathrm dt.
\]
The last equality is the nonnegative layer-cake identity, including
possible $T=+\infty$. All finite cutoffs are bounded, so $L_N\leq2N$.
Define $\Psi:[0,2]\to[0,\infty)$ by
\[
 \Psi(x)=\sup_{0\leq a\leq1}(ax-a^3/3)
 =\begin{cases}(2/3)x^{3/2},&0\leq x\leq1,\\
                x-1/3,&1\leq x\leq2.
   \end{cases}
\]
These restricted duals are not definitions of full admissibility.

\begin{theorem}\label{thm:duals}
The excess single-pencil dual is
\[
 C_{\rm pen}=\sup_{N\geq1,\,0\leq a\leq1}
       \bigl(E(u_0+aD_N)-\int(u_0+aD_N)\,\mathrm d\mu-J_0\bigr)
          =\sup_{N\geq1}N\Psi(L_N/N).
\]
If $B=\sup_{N\geq1} L_N/N^{1/3}$, then, in the extended sense,
\[
 \frac5{6\sqrt2}B^{3/2}\leq C_{\rm pen}\leq\frac23B^{3/2}.
\]
It is finite exactly when
$\sup_{t\geq1}t^{2/3}\mu\{T\geq t\}<\infty$.

The excess dual on all finite combinations in Proposition~\ref{prop:energy}
is instead
\[
 C_{\rm rad}=\Psi(L_1)+\int_1^\infty\Psi(F_\mu(t))\,\mathrm dt.
\]
In particular it is finite exactly when
\[
                    \int_1^\infty F_\mu(t)^{3/2}\,\mathrm dt<\infty.
\]
If this integral diverges, the bounded-class dual
$E^*_{\C}(\mu)=\sup_{u\in\C}(E(u)-\int u\,\mathrm d\mu)$
is infinite, under the standing assumptions. Finiteness of either
restricted dual does not imply $E^*_{\C}(\mu)<\infty$.
\end{theorem}
\begin{proof}
For a single cutoff Proposition~\ref{prop:energy} gives the excess
$aL_N-a^3N/3$, whose exact maximum is $N\Psi(L_N/N)$.
For $0\leq x\leq2$ direct differentiation shows
\[
 \frac5{6\sqrt2}x^{3/2}\leq\Psi(x)\leq\frac23x^{3/2}.
\]
On $[0,1]$ the upper bound is equality. On $[1,2]$ the ratio
$(x-1/3)/x^{3/2}$ decreases from $2/3$ to $5/(6\sqrt2)$.
Taking suprema proves the claimed bounds.
If $Q=\sup_{t\geq1}t^{2/3}\mu\{T\geq t\}$, then
$t\mu\{T\geq t\}\leq L_t$ gives $Q\leq B$.
Conversely $L_N\leq2+Q\int_1^N t^{-2/3}\,\mathrm dt$
for $N\geq1$, so $B\leq2+3Q$. This proves the pencil equivalence.

For every finite combination, layer cake and the proposition give excess
\[
       \int_0^\infty\bigl(h(t)F_\mu(t)-h(t)^3/3\bigr)\,\mathrm dt.
\]
On $[0,1)$, $h$ is constant, making this contribution
$h(0)L_1-h(0)^3/3\leq\Psi(L_1)$.
On $[1,\infty)$ pointwise maximization gives the displayed integral
upper bound. This upper bound is the supremum: choose
\[
 h_*(t)=\begin{cases}
    \min(1,\sqrt{L_1}),&0\leq t<1,\\
    \min(1,\sqrt{F_\mu(t)}),&t\geq1.
 \end{cases}
\]
Since $F_\mu$ is nonincreasing and $F_\mu(1)\leq L_1$, this profile
is nonincreasing. Choose bounded intervals $[1,m]$, sample the tail
profile at right endpoints of successively refined dyadic partitions,
and set it to zero beyond $m$. Keep its stated constant on $[0,1)$.
The resulting compact finite step profiles have only downward jumps,
at cutoffs at least one, and increase almost everywhere to $h_*$ along
a nested choice of domains and partitions. They correspond to actual
finite convex combinations of the $D_N$. Their excess integrands
increase to the optimal nonnegative integrand, because
$a\mapsto aF-a^3/3$ increases for $a\leq\min(1,\sqrt F)$.
Monotone convergence proves equality even if the supremum is infinite.
The two-sided bounds for $\Psi$ prove the last integral criterion.
Finally all these competitors belong to $\C$, so an infinite restricted
supremum forces the full bounded supremum to be infinite.
\end{proof}

\section{A sharp analytic atomic example and its limitation}
For $\beta>0$ define an actual analytic measure
\[
 \mu_\beta=2\sum_{n=1}^\infty
             \bigl(n^{-\beta}-(n+1)^{-\beta}\bigr)\delta_{\xi_n}.
\]
The coefficients are positive and telescope to raw mass two. This is a
Radon measure: its finite atomic partial sums converge in total variation,
and the remaining tail mass tends to zero on the compact Hausdorff carrier.
It has no rigid atom. Since $T(\xi_n)=n$,
\[
 \mu_\beta\{T\geq t\}=2\lceil t\rceil^{-\beta}\quad(t\geq1),
 \qquad F_{\mu_\beta}(t)=2(n+1)^{-\beta}\quad(n<t<n+1).
\]
Theorem~\ref{thm:duals} shows
\[
 C_{\rm pen}(\mu_\beta)<\infty\quad\Longleftrightarrow\quad\beta\geq2/3,
 \qquad
 C_{\rm rad}(\mu_\beta)<\infty\quad\Longleftrightarrow\quad\beta>2/3.
\]
For the latter assertion the large-$n$ summands are exactly
$(2/3)2^{3/2}(n+1)^{-3\beta/2}$; the elementary power-series criterion
gives the stated strict threshold. At $\beta=2/3$ the full bounded dual
is therefore infinite although the entire single-pencil dual is finite.
This rules out the weak $t^{-2/3}$ tail estimate as a sufficient
bounded-class energy criterion.

For comparison, a fixed $\mu=2\delta_{\xi_{N_0}}$, $N_0\geq1$, has
$L_N=2\min(N,N_0)$ and
\[
 J_0=1/6,\qquad C_{\rm pen}=C_{\rm rad}=5N_0/3.
\]
The pencil maximum is attained at $N=N_0,a=1$; for the radial dual
$F=2$ on $(0,N_0)$ and zero thereafter. Thus the two restricted duals
both equal $1/6+5N_0/3$ after restoring $J_0$.
Neither calculation proves finiteness of the full bounded dual for this
measure, existence of an equation solution, or any universal upper bound
on the atom of a Monge--Amp\`ere measure.

\section{Proof and source boundaries}
This is a new complete conditional multiscale argument. The full original
P16.7 remains open: it asks for an intrinsic unbounded finite-energy
operator, full mass and decreasing continuity, useful separated weak
compactness, all analytic-obstacle envelopes and orthogonality, full
energy variation, comparison, and existence and uniqueness for every
admissible measure. None is supplied by restricted dual calculations.
In particular the critical measure above is excluded by the stronger
bounded dual, so it is not an admissible counterexample to that equation.

The geometry and bounded calculus are exact assumptions from the complete
unreviewed five-note continuation 383, especially note~3. Its full proof,
distinct actual and expanded notes, and scope were read. No external
review of these inputs is claimed. CLD-v2, 25 July 2025,
Sections~5.8.1--6 are the directly applied integration statements and
immediate proofs for strong support, calibration and Radon measures.
Their entire retained text at LF~6720--6840 was freshly read for this paper.
On the two translated identity moments the saturated lattice is
$\mathbb Z^2$ and the degree is one. Full deeper integration foundations,
the source's 608 pages, and recursive proof trees were not newly read or
proved. The gradient, finite-step integration, optimization, tail and
atomic-measure arguments above are proved here; no supplied or native
program was run or imported. All full primary bodies and literal units
remain local. Compiler success and paper acceptance are separate pending
checks. The already open P16.4 editor remains unchanged.
\end{document}
